
\subsubsection{Distributed Packing and Padding}\label{sec:train_infra_pnp}



Due to the integrity constraints of video data and the variability in video lengths and resolutions, we adopt a Packing and Padding (PnP) strategy~\citep{sirluk2024efficient, kundu2024enhancingtrainingefficiencyusing} to batch video samples in a way that minimizes excessive padding and reduces unnecessary computational overhead in distributed training scenarios. Moreover, the data composition is frequently adjusted during the training of \magi (See Sec.~\ref{sec:data_adjustment}), and to accommodate such flexibility, we employ an online PnP strategy instead of a offline approach.

The core idea of PnP is to efficiently utilize GPU resources by concatenating multiple short sequences into a batch while minimizing redundant filling. The offline formulation of this problem aligns with the classic bin-packing problem: given a set of input samples, the goal is to pack them into a set of bins, each with a fixed capacity \texttt{max\_length}, while minimizing overall unused space. Although this problem is NP-complete, it can be efficiently approximated in practice using the First-Fit Decreasing (FFD)~\citep{dosa2007tight} greedy algorithm.

In our online setting, we must process streaming data inputs while ensuring compatibility with the 3D parallelism strategy employed during training. To this end, we reformulate the problem as follows: given M candidate samples, we aim to pack them into N bins of size \texttt{max\_length}, minimizing overall space waste. Here, M denotes the size of the candidate pool with M $\gg$ N; N must be divisible by the \texttt{DP\_SIZE}; and \texttt{max\_length} must be divisible by \texttt{TP\_SIZE}$\times$\texttt{CP\_SIZE}.

In practice, we extend the FFD algorithm with custom heuristics to support efficient online packing under these constraints. This approach enables us to achieve a $99\%$ capacity utilization rate and the differences between different DP groups can be neglected, thus substantially reducing computational overhead during training.
































